% ============================================================ Capítulo 6
\chapter{Arquitectura de software}\label{cap:arquitectura}

\section{Estilo arquitectónico}
Se adopta una \textbf{arquitectura por capas con un núcleo de servicios compartido} (estilo hexagonal
simplificado): las interfaces ---Streamlit, API, CLI y MCP--- son adaptadores de entrada que invocan una única
fachada, \texttt{ServicioAsistente}; los proveedores de LLM, el almacenamiento de archivos y la base relacional
son adaptadores de salida. La dirección de dependencias se verifica automáticamente con \texttt{scripts/verificar\_capas.py} (capítulo~\ref{cap:pruebas}).

\section{Diagrama de contexto}
\begin{figure}[H]\centering
\begin{tikzpicture}[node distance=1.2cm and 1.6cm]
  \node[clave=4.2cm, minimum height=1.6cm] (sis) {\textbf{Sistema de conocimiento Fanalca}\\Base de conocimiento + Resumen, FAQ y Q\&A};
  \node[actor, left=2.2cm of sis, yshift=1.1cm] (cli) {Cliente};
  \node[actor, left=2.2cm of sis, yshift=-1.3cm] (dev) {Equipo de\\desarrollo};
  \node[externo=2.9cm, right=2.0cm of sis, yshift=1.4cm] (web) {Sitio oficial\\\texttt{fanalca.com}};
  \node[externo=2.9cm, right=2.0cm of sis, yshift=0cm] (llm) {Proveedor LLM\\Ollama / OpenAI / Anthropic};
  \node[externo=2.9cm, right=2.0cm of sis, yshift=-1.4cm] (gh) {GitHub\\(repositorio y CI)};
  \node[externo=2.9cm, below=1.3cm of sis] (ag) {Consolas agénticas\\OpenCode / Claude Code};
  \draw[flecha] (cli) -- node[etq,above,sloped]{preguntas} (sis);
  \draw[flecha] (dev) -- node[etq,below,sloped]{opera y evalúa} (sis);
  \draw[flecha] (sis) -- node[etq,above,sloped]{HTTPS + robots.txt} (web);
  \draw[flecha] (sis) -- node[etq,above]{HTTP} (llm);
  \draw[flecha] (dev.south east) -- node[etq,below,sloped]{push / PR} (gh);
  \draw[flecha] (ag) -- node[etq,right]{MCP (stdio)} (sis);
\end{tikzpicture}
\caption{Diagrama de contexto (C4 nivel 1)}\label{fig:contexto}
\end{figure}
El sistema tiene dos actores humanos y cuatro sistemas externos. La única comunicación saliente hacia Internet
del pipeline es el rastreo del sitio oficial; la generación puede ser totalmente local con Ollama.

\section{Arquitectura general y contenedores}
\begin{figure}[H]\centering
\begin{tikzpicture}[node distance=0.55cm and 0.55cm]
  \node[caja=2.3cm] (ui) {Streamlit\\\scriptsize frontend/app.py};
  \node[caja=2.3cm, right=of ui] (api) {FastAPI\\\scriptsize api/main.py};
  \node[caja=2.3cm, right=of api] (cli) {CLI \texttt{fabrica}\\\scriptsize cli.py};
  \node[caja=2.3cm, right=of cli] (mcp) {Servidor MCP\\\scriptsize mcp\_server/};
  \node[clave=10.4cm, below=1.0cm of $(api)!0.5!(cli)$] (svc) {\textbf{ServicioAsistente} (services/) --- responder · generar\_faq · resumir · construir · rastrear · evaluar};
  \node[caja=2.6cm, below=1.0cm of svc.south west, anchor=north west] (llm) {Cadenas LCEL\\\scriptsize llm/ (prompts v1--v3,\\salvaguardas)};
  \node[caja=2.6cm, right=0.5cm of llm] (kb) {BaseConocimiento\\\scriptsize knowledge/ (contexto,\\BM25)};
  \node[caja=2.6cm, right=0.5cm of kb] (pipe) {Pipeline Medallion\\\scriptsize pipelines/ + scraping/};
  \node[almacen, below=0.9cm of kb] (gold) {Gold\\\scriptsize JSONL + TXT};
  \node[almacen, below=0.9cm of pipe] (bronze) {Bronze/Silver\\\scriptsize HTML + JSONL};
  \node[almacen, right=0.6cm of pipe] (db) {SQLite\\\scriptsize auditoría};
  \node[externo=2.4cm, below=0.9cm of llm] (prov) {Ollama / API};
  \node[externo=2.2cm, right=1.6cm of bronze] (web) {fanalca.com};
  \foreach \n in {ui,api,cli,mcp} \draw[flecha] (\n) -- (\n |- svc.north);
  \draw[flecha] (svc.south -| llm) -- (llm);
  \draw[flecha] (svc.south -| kb) -- (kb);
  \draw[flecha] (svc.south -| pipe) -- (pipe);
  \draw[flecha] (svc.east) -| (db);
  \draw[flecha] (llm) -- (prov);
  \draw[flecha] (kb) -- (gold);
  \draw[flecha] (pipe) -- (bronze);
  \draw[flecha] (pipe.south west) -- (gold.north east);
  \draw[flechad] (bronze) -- node[etq,below]{HTTPS} (web);
  \draw[flechad] (mcp.south east) to[out=-60,in=0] node[etq,right,pos=0.7]{solo lectura} (gold.east);
\end{tikzpicture}
\caption{Arquitectura general y contenedores (C4 nivel 2)}\label{fig:arquitectura}
\end{figure}

\section{Diagrama de componentes}
\begin{figure}[H]\centering
\begin{tikzpicture}[node distance=0.35cm]
  \tikzset{mod/.style={caja=3.1cm,minimum height=0.7cm}}
  \node[mod] (c1) {core: config, registro, excepciones};
  \node[mod, right=of c1] (c2) {domain: modelos, puertos};
  \node[mod, above=0.9cm of c1] (s1) {scraping: robots, cliente\_http, sitemap, extracción, rastreador};
  \node[mod, right=of s1] (p1) {pipelines: bronze, silver, fragmentación, gold, calidad};
  \node[mod, right=of p1] (k1) {knowledge: base, recuperación BM25};
  \node[mod, above=0.9cm of p1] (l1) {llm: catálogo prompts, cadenas, proveedores, salvaguardas, extractivo};
  \node[mod, right=of l1] (r1) {repositories: orm, auditoría};
  \node[mod, above=0.9cm of l1] (sv) {services: asistente, evaluación};
  \node[mod, left=of sv] (if1) {api · cli};
  \node[mod, right=of sv] (if2) {mcp\_server · frontend};
  \draw[flecha] (if1) -- (sv); \draw[flecha] (if2) -- (sv);
  \draw[flecha] (sv) -- (l1); \draw[flecha] (sv) -- (r1); \draw[flecha] (sv.south west) -- (s1.north);
  \draw[flecha] (l1) -- (k1); \draw[flecha] (k1) -- (p1); \draw[flecha] (p1) -- (s1);
  \draw[flecha] (s1) -- (c1); \draw[flecha] (s1) -- (c2);
  \node[font=\sffamily\scriptsize, align=left, right=0.4cm of c2] {Nivel 0\\(sin dependencias\\internas)};
  \node[font=\sffamily\scriptsize, align=left, right=0.4cm of if2] {Nivel 6};
\end{tikzpicture}
\caption{Componentes por capa; las flechas indican dependencia permitida}\label{fig:componentes}
\end{figure}

\section{Secuencia de una consulta Q\&A}
\begin{figure}[H]\centering
\begin{tikzpicture}[x=2.05cm,y=-0.62cm]
  \foreach \x/\n in {0/Usuario,1/Streamlit,2/Servicio,3/Salvaguardas,4/Base,5/LLM,6/Auditoría}
    {\node[caja=1.7cm,minimum height=0.6cm,fill=fnaclaro] at (\x,0) {\n};
     \draw[dashed,fnagris!60] (\x,0.6) -- (\x,15.6);}
  \newcommand{\msg}[4]{\draw[flecha] (#1,#3) -- node[etq,above]{\scriptsize #4} (#2,#3);}
  \newcommand{\ret}[4]{\draw[flechad] (#1,#3) -- node[etq,above]{\scriptsize #4} (#2,#3);}
  \msg{0}{1}{1.5}{pregunta}
  \msg{1}{2}{2.5}{responder()}
  \msg{2}{3}{3.5}{detectar\_inyeccion}
  \ret{3}{2}{4.5}{válida}
  \msg{2}{4}{5.5}{armar\_contexto(modo)}
  \ret{4}{2}{6.5}{<contexto> fragmentos}
  \msg{2}{5}{7.5}{prompt v3}
  \ret{5}{2}{8.5}{JSON respuesta+fuentes}
  \msg{2}{3}{9.5}{validar citas y fundamentación}
  \ret{3}{2}{10.5}{fundamentada}
  \msg{2}{6}{11.5}{registrar\_interaccion}
  \ret{2}{1}{12.8}{RespuestaQA}
  \ret{1}{0}{14.0}{respuesta + fuentes}
  \node[draw=fnarojo,fill=fnarojo!5,font=\sffamily\scriptsize,text width=3.2cm,align=left] at (4.6,3.1)
    {Si hay inyección: respuesta de bloqueo, sin invocar el LLM};
\end{tikzpicture}
\caption{Diagrama de secuencia de la consulta Q\&A}\label{fig:secuencia}
\end{figure}

\section{Diagrama de actividades del pipeline}
\begin{figure}[H]\centering
\begin{tikzpicture}[node distance=0.42cm]
  \tikzset{act/.style={caja=4.6cm,rounded corners=8pt,minimum height=0.6cm}, dec/.style={diamond,draw=fnagris,aspect=2.4,font=\sffamily\scriptsize,inner sep=1pt,fill=oro!10}}
  \node[circle,fill=fnagris,minimum size=0.3cm] (ini) {};
  \node[act, below=of ini] (a1) {Cargar robots.txt y sitemaps};
  \node[act, below=of a1] (a2) {Tomar URL de la cola (BFS)};
  \node[dec, below=of a2] (d1) {¿Permitida?};
  \node[act, below=of d1] (a3) {GET con pausa, timeout y reintentos};
  \node[dec, below=of a3] (d2) {¿200 + HTML nuevo?};
  \node[act, below=of d2] (a4) {Guardar en Bronze + encolar enlaces};
  \node[dec, below=of a4] (d3) {¿Cola vacía o límite?};
  \node[act, below=of d3] (a5) {Silver: extraer, limpiar, deduplicar};
  \node[act, below=of a5] (a6) {Gold: Markdown, fragmentos, consolidado};
  \node[act, below=of a6] (a7) {Validar calidad};
  \node[circle,draw=fnagris,double,fill=fnagris,minimum size=0.3cm, below=of a7] (fin) {};
  \node[act, right=1.6cm of d1, text width=2.8cm] (om) {Registrar omisión};
  \node[act, right=1.6cm of d2, text width=2.8cm] (er) {Registrar error/duplicado};
  \foreach \a/\b in {ini/a1,a1/a2,a2/d1,a3/d2,a4/d3,a5/a6,a6/a7,a7/fin} \draw[flecha] (\a) -- (\b);
  \draw[flecha] (d1) -- node[etq,right]{sí} (a3);
  \draw[flecha] (d2) -- node[etq,right]{sí} (a4);
  \draw[flecha] (d3) -- node[etq,right]{sí} (a5);
  \draw[flecha] (d1) -- node[etq,above]{no} (om);
  \draw[flecha] (d2) -- node[etq,above]{no} (er);
  \draw[flecha] (om.north) |- (a2.east);
  \draw[flecha] (er.east) -- ++(0.5,0) |- (a2.east);
  \draw[flecha] (d3.west) -- node[etq,above]{no} ++(-1.4,0) |- (a2.west);
\end{tikzpicture}
\caption{Diagrama de actividades: rastreo y construcción de la base}\label{fig:actividades}
\end{figure}

\section{Modelo de dominio}
\begin{figure}[H]\centering
\begin{tikzpicture}[node distance=0.9cm and 1.1cm]
  \tikzset{clase/.style={rectangle split,rectangle split parts=2,draw=fnagris,fill=white,
    font=\sffamily\scriptsize,text width=3.6cm,rectangle split part fill={fnaazul!10,white}}}
  \node[clase] (rc) {\textbf{RegistroCaptura}\nodepart{second}url, url\_final, fecha\_captura, estado\_http, sha256, archivo\_html};
  \node[clase, right=of rc] (dl) {\textbf{DocumentoLimpio}\nodepart{second}id, url, titulo, seccion, idioma, texto, sha256\_texto, origen};
  \node[clase, right=of dl] (fr) {\textbf{Fragmento}\nodepart{second}id (doc\#NNN), documento\_id, url, seccion, encabezado, texto, tokens\_aprox, origen};
  \node[clase, below=1.4cm of fr] (fu) {\textbf{Fuente}\nodepart{second}fragmento\_id, titulo, url};
  \node[clase, left=of fu] (qa) {\textbf{RespuestaQA}\nodepart{second}pregunta, respuesta, fundamentada, confianza, version\_prompt, modo\_contexto, latencia\_ms};
  \node[clase, left=of qa] (it) {\textbf{Interaccion} (tabla)\nodepart{second}id, fecha, tarea, canal, entrada, salida, proveedor, modelo, fuentes};
  \draw[flecha] (rc) -- node[etq,above]{1 : 0..1} (dl);
  \draw[flecha] (dl) -- node[etq,above]{1 : 1..*} (fr);
  \draw[flecha] (qa) -- node[etq,above]{1 : 0..*} (fu);
  \draw[flecha] (fu) -- node[etq,right]{referencia} (fr);
  \draw[flechad] (qa) -- node[etq,above]{se audita} (it);
\end{tikzpicture}
\caption{Modelo de dominio}\label{fig:dominio}
\end{figure}

\section{Registro de decisiones de arquitectura}
\begin{table}[H]\centering\small
\caption{Resumen de ADR (texto completo en \texttt{docs/arquitectura/adr/})}\label{tab:adr}
\begin{tabularx}{\linewidth}{@{}l X X@{}}\toprule
ADR & Decisión & Justificación principal\\\midrule
001 & Python 3.12 con uv & Estándar del curso; reproducibilidad con \texttt{uv.lock}.\\
002 & Medallion en archivos JSONL/Markdown & Corpus pequeño, revisable y reprocesable sin red.\\
003 & SQLite + SQLAlchemy + Alembic para auditoría & Sin servidor; migrable a PostgreSQL cambiando la URL.\\
004 & LangChain LCEL; Ollama por defecto & Recomendación del taller; privacidad y costo cero por token.\\
005 & Contexto completo con respaldo BM25 & Cumple el taller y evita truncamiento con modelos locales.\\
006 & Streamlit & Coherencia con la Clase 1; pruebas con AppTest.\\
007 & FastAPI como extensión & Integración programática; contrato OpenAPI.\\
008 & Servidor MCP propio de solo lectura (SDK 2.x) & Acceso de agentes con mínimo privilegio.\\
009 & Prompts versionados en YAML & Experimentación trazable y comparable.\\
010 & Redes sociales solo como consulta manual & Términos de uso y derechos de autor.\\\bottomrule
\end{tabularx}
\end{table}

% ============================================================ Capítulo 7
\chapter{Pila tecnológica}\label{cap:pila}
\begin{longtable}{@{}P{2.8cm}P{3.4cm}P{9.0cm}@{}}
\caption{Selección tecnológica justificada (versiones resueltas en \texttt{uv.lock} durante la construcción)}\label{tab:pila}\\\toprule
Componente & Selección & Justificación frente a los requisitos\\\midrule\endfirsthead
\toprule Componente & Selección & Justificación\\\midrule\endhead
Sistema operativo & Windows 11 & Entorno principal exigido; scripts PowerShell y CI en \texttt{windows-latest}.\\
Terminal & Windows Terminal + PowerShell & Incluidos en Windows 11; los scripts se verificaron sintácticamente con PowerShell 7.4.6.\\
Control de versiones & Git + GitHub + GitHub CLI & Entrega en GitHub exigida por la rúbrica; \texttt{gh} simplifica autenticación y creación del repositorio.\\
Editor & Visual Studio Code & Extensiones de Python, Ruff, YAML, TypeScript y LaTeX Workshop.\\
Gestión Python & uv (0.11.32 en construcción), Python 3.12.13 & Obligatorio por la Clase 2.\\
Consolas agénticas & OpenCode y Claude Code & Exigidas por la Clase 3 y el enunciado; configuraciones versionadas.\\
Integración de herramientas & MCP SDK 2.3.0 & Servidor propio; estándar abierto soportado por ambas consolas.\\
Scraping & requests 2.34, BeautifulSoup 4.15 + lxml 6.1, trafilatura 2.3; Playwright opcional & El taller nombra requests, BeautifulSoup y Selenium; trafilatura (Clase 2) mejora la extracción del cuerpo; Playwright sustituye a Selenium para páginas dinámicas. Scrapy no se incorpora: el volumen no lo justifica.\\
Orquestación LLM & LangChain core 1.7, langchain-ollama 1.1, langchain-openai 1.7, langchain-anthropic (extra) & Recomendado por el taller; LCEL permite cambiar de modelo sin reescribir cadenas.\\
Modelo por defecto & Ollama \texttt{llama3.2:3b} & Ejemplo del taller y de la Clase 1; ejecución local.\\
Backend & FastAPI 0.162 + Uvicorn & Validación con Pydantic; OpenAPI automático.\\
Validación & Pydantic 2.14 + pydantic-settings & Modelos de dominio y configuración por entorno.\\
Frontend & Streamlit 1.66 & Requisito del taller (Streamlit/Gradio).\\
Persistencia & SQLite + SQLAlchemy 2.1 + Alembic 1.20; PostgreSQL 17 (psycopg) & SQLite sin servidor para desarrollo; PostgreSQL para el despliegue de la versión 2.0 con las mismas migraciones (ADR-003, ADR-012). MongoDB descartado.\\
Procesamiento & pandas 3.1, matplotlib & Tablas en Streamlit y gráficas por etapa del pipeline. Polars descartado: sin beneficio a esta escala. Sin notebooks ni Anaconda (ADR-014).\\
Geometría 3D (v2.0) & trimesh + manifold3d & Generación procedural por CLI y exportación GLB; Blender CLI y CadQuery evaluados y descartados (ADR-011).\\
Visor 3D (v2.0) & Vite, React 19, TypeScript, three.js, React Three Fiber & Escena WebGL con selección, sección y animación por fotograma; servido por FastAPI en \texttt{/app} (ADR-013).\\
Pruebas & pytest 9, pytest-asyncio, pytest-cov, responses, httpx, pgserver; Vitest y Playwright & Pruebas sin red ni GPU; PostgreSQL embebido; visor probado con Chrome.\\
Calidad & Ruff 0.17, mypy 2.4, pre-commit & Estilo, tipos y ganchos antes del commit.\\
API & OpenAPI/Swagger, Postman & Contrato exportado y colección con pruebas.\\
Contenedores & Docker Desktop (opcional) & Despliegue repetible de Ollama, API e interfaz.\\
Documentación & LaTeX + TikZ; Mermaid y PlantUML como fuentes & Manual compilable con diagramas vectoriales.\\
Integración continua & GitHub Actions & Matriz Ubuntu/Windows con \texttt{uv sync --locked}.\\\bottomrule
\end{longtable}
